\documentclass[12pt]{article}
\usepackage[margin=2.5cm]{geometry}
\usepackage{array}
\usepackage{hyperref}
\usepackage{longtable}
\usepackage{graphicx}

\title{Informe: Implementacion Automatizada de una API REST de Gestion de Pedidos en AWS}
\author{Grupo de trabajo}
\date{\today}

\begin{document}
\maketitle

\section{Resumen}

Se desarrollo una API REST para una plataforma basica de gestion de pedidos. La solucion permite crear pedidos, listar pedidos con filtros y ordenamiento, consultar pedidos por identificador y actualizar su estado. Ademas, se definio una automatizacion DevOps en AWS usando control de versiones, integracion continua, generacion de artefactos y despliegue automatico hacia un entorno de pruebas.

\section{Requisitos de negocio}

\begin{longtable}{|p{5cm}|p{9cm}|}
\hline
\textbf{Requisito} & \textbf{Cumplimiento} \\
\hline
Crear pedidos & Se implemento el endpoint \texttt{POST /orders}, que recibe cliente, productos, cantidades y direccion de envio. El pedido queda con estado inicial \texttt{Pendiente}. \\
\hline
Visualizar pedidos & Se implemento \texttt{GET /orders}. Permite filtrar por estado y ordenar por fecha de creacion, estado o cliente. \\
\hline
Actualizar estado & Se implemento \texttt{PATCH /orders/\{id\}/status}. Los estados permitidos son \texttt{Pendiente}, \texttt{En Proceso}, \texttt{Enviado} y \texttt{Entregado}. \\
\hline
\end{longtable}

\section{Arquitectura de la aplicacion}

La API fue desarrollada con Node.js usando el modulo HTTP nativo, sin dependencias externas obligatorias. Esta decision permite ejecutar pruebas y builds de forma simple en AWS CodeBuild. Para fines academicos, los pedidos se almacenan en memoria; en un escenario productivo se recomienda incorporar Amazon DynamoDB o Amazon RDS.

\subsection{Endpoints principales}

\begin{itemize}
  \item \texttt{GET /health}: verifica disponibilidad del servicio.
  \item \texttt{POST /orders}: crea un pedido.
  \item \texttt{GET /orders}: lista pedidos con filtros y ordenamiento.
  \item \texttt{GET /orders/\{id\}}: consulta un pedido especifico.
  \item \texttt{PATCH /orders/\{id\}/status}: actualiza el estado del pedido.
\end{itemize}

\section{Automatizacion DevOps en AWS}

\subsection{Control de versiones}

El codigo fuente y la infraestructura como codigo se gestionan en un repositorio AWS CodeCommit. La plantilla de CloudFormation crea el repositorio base para el proyecto.

\textbf{Evidencia sugerida:} insertar captura del repositorio CodeCommit con los archivos \texttt{src/}, \texttt{buildspec.yml}, \texttt{Dockerfile} e \texttt{infrastructure/cloudformation-pipeline.yml}.

\subsection{Integracion continua}

AWS CodePipeline se configura para iniciar el flujo cuando se confirman cambios en la rama principal. AWS CodeBuild usa el archivo \texttt{buildspec.yml} para instalar dependencias, ejecutar pruebas unitarias y preparar el artefacto de despliegue.

\textbf{Evidencia sugerida:} insertar captura de CodeBuild mostrando una ejecucion exitosa y el log con \texttt{npm test}.

\subsection{Entrega continua}

El pipeline incluye una etapa de despliegue hacia AWS Elastic Beanstalk. El artefacto generado por CodeBuild contiene el codigo de la API y el \texttt{Procfile}, por lo que Elastic Beanstalk puede iniciar el servidor con Node.js.

\textbf{Evidencia sugerida:} insertar captura de CodePipeline con las etapas Source, Build y DeployTest en verde.

\subsection{Monitoreo}

Elastic Beanstalk publica eventos y metricas en Amazon CloudWatch. Esto permite revisar logs, errores y estado del entorno desplegado.

\textbf{Evidencia sugerida:} insertar captura de eventos de Elastic Beanstalk o logs en CloudWatch.

\section{Infraestructura como codigo}

Se creo la plantilla \texttt{infrastructure/cloudformation-pipeline.yml}, que provisiona:

\begin{itemize}
  \item Repositorio AWS CodeCommit.
  \item Bucket S3 para artefactos.
  \item Proyecto AWS CodeBuild.
  \item Pipeline AWS CodePipeline.
  \item Aplicacion y entorno AWS Elastic Beanstalk.
  \item Roles IAM necesarios para CodeBuild y CodePipeline.
\end{itemize}

\section{Pruebas}

Las pruebas unitarias verifican:

\begin{itemize}
  \item Creacion de pedidos validos.
  \item Listado y consulta de pedidos.
  \item Actualizacion del estado.
  \item Rechazo de datos incompletos o estados invalidos.
\end{itemize}

\section{Conclusiones}

La solucion cumple los requisitos funcionales del caso de estudio y define un flujo DevOps automatizado en AWS. El pipeline permite que cada cambio confirmado en el repositorio sea construido, probado y desplegado automaticamente. Como mejora futura, se recomienda agregar persistencia con DynamoDB o RDS, autenticacion de usuarios y despliegue separado para ambientes de pruebas y produccion.

\end{document}
